\documentclass{article}
\usepackage[utf8]{inputenc}
\usepackage[spanish]{babel}
\usepackage{amsmath}
\usepackage{amsfonts}

\begin{document}
	
	\section*{Resolución de Sistema Lineal - Paso a Paso}
	
	\subsection*{1. Construcción de Matrices}
	Dado el sistema:
	\[
	\begin{cases}
		x + y + z = 6 \\
		2x - y + z = 3 \\
		3x + 2y - 2z = 2
	\end{cases}
	\]
	
	Definimos las matrices $A$, $X$ y $B$:
	\[
	A = \begin{pmatrix} 1 & 1 & 1 \\ 2 & -1 & 1 \\ 3 & 2 & -2 \end{pmatrix}, \quad 
	X = \begin{pmatrix} x \\ y \\ z \end{pmatrix}, \quad 
	B = \begin{pmatrix} 6 \\ 3 \\ 2 \end{pmatrix}
	\]
	
	\subsection*{2. Cálculo del Determinante $\det(A)$}
	Utilizando el método de la matriz ampliada (Sarrus):
	\[
	\det(A) = \begin{vmatrix} 1 & 1 & 1 \\ 2 & -1 & 1 \\ 3 & 2 & -2 \end{vmatrix}
	\]
	Cálculo:
	\[
	\det(A) = [(1)(-1)(-2) + (1)(1)(3) + (1)(2)(2)] - [(3)(-1)(1) + (2)(1)(1) + (-2)(2)(1)]
	\]
	\[
	\det(A) = [2 + 3 + 4] - [-3 + 2 - 4] = 9 - (-5) = \mathbf{14}
	\]
	Como $\det(A) = 14 \neq 0$, la matriz es invertible.
	
	\subsection*{3. Cálculo de $A^{-1}$ por Gauss-Jordan}
	Ampliamos con la identidad $[A | I]$:
	\[
	\left( \begin{array}{ccc|ccc} 1 & 1 & 1 & 1 & 0 & 0 \\ 2 & -1 & 1 & 0 & 1 & 0 \\ 3 & 2 & -2 & 0 & 0 & 1 \end{array} \right)
	\]
	
	\textbf{Paso 1: Ceros en la primera columna}
	$R_2 \leftarrow R_2 - 2R_1$ y $R_3 \leftarrow R_3 - 3R_1$:
	\[
	\left( \begin{array}{ccc|ccc} 1 & 1 & 1 & 1 & 0 & 0 \\ 0 & -3 & -1 & -2 & 1 & 0 \\ 0 & -1 & -5 & -3 & 0 & 1 \end{array} \right)
	\]
	
	\textbf{Paso 2: Intercambio para pivote unitario}
	$R_2 \leftrightarrow R_3$:
	\[
	\left( \begin{array}{ccc|ccc} 1 & 1 & 1 & 1 & 0 & 0 \\ 0 & -1 & -5 & -3 & 0 & 1 \\ 0 & -3 & -1 & -2 & 1 & 0 \end{array} \right)
	\]
	Multiplicamos $R_2$ por $-1$ ($R_2 \leftarrow -R_2$):
	\[
	\left( \begin{array}{ccc|ccc} 1 & 1 & 1 & 1 & 0 & 0 \\ 0 & 1 & 5 & 3 & 0 & -1 \\ 0 & -3 & -1 & -2 & 1 & 0 \end{array} \right)
	\]
	
	\textbf{Paso 3: Ceros en la segunda columna}
	$R_1 \leftarrow R_1 - R_2$ y $R_3 \leftarrow R_3 + 3R_2$:
	\[
	\left( \begin{array}{ccc|ccc} 1 & 0 & -4 & -2 & 0 & 1 \\ 0 & 1 & 5 & 3 & 0 & -1 \\ 0 & 0 & 14 & 7 & 1 & -3 \end{array} \right)
	\]
	
	\textbf{Paso 4: Normalizar tercera fila}
	$R_3 \leftarrow \frac{1}{14}R_3$:
	\[
	\left( \begin{array}{ccc|ccc} 1 & 0 & -4 & -2 & 0 & 1 \\ 0 & 1 & 5 & 3 & 0 & -1 \\ 0 & 0 & 1 & 1/2 & 1/14 & -3/14 \end{array} \right)
	\]
	
	\textbf{Paso 5: Ceros en la tercera columna (Final)}
	$R_1 \leftarrow R_1 + 4R_3$ y $R_2 \leftarrow R_2 - 5R_3$:
	\[
	\left( \begin{array}{ccc|ccc} 1 & 0 & 0 & 0 & 4/14 & 2/14 \\ 0 & 1 & 0 & 7/14 & -5/14 & 1/14 \\ 0 & 0 & 1 & 7/14 & 1/14 & -3/14 \end{array} \right)
	\]
	
	Extrayendo $A^{-1}$:
	\[
	A^{-1} = \frac{1}{14} \begin{pmatrix} 0 & 4 & 2 \\ 7 & -5 & 1 \\ 7 & 1 & -3 \end{pmatrix}
	\]
	
	\subsection*{4. Solución Final $X = A^{-1}B$}
	\[
	\begin{pmatrix} x \\ y \\ z \end{pmatrix} = \frac{1}{14} \begin{pmatrix} 0 & 4 & 2 \\ 7 & -5 & 1 \\ 7 & 1 & -3 \end{pmatrix} \begin{pmatrix} 6 \\ 3 \\ 2 \end{pmatrix}
	\]
	Multiplicación:
	\begin{itemize}
		\item $x = \frac{1}{14} [0(6) + 4(3) + 2(2)] = \frac{16}{14} = \mathbf{\frac{8}{7}}$
		\item $y = \frac{1}{14} [7(6) - 5(3) + 1(2)] = \frac{29}{14} = \mathbf{\frac{29}{14}}$
		\item $z = \frac{1}{14} [7(6) + 1(3) - 3(2)] = \frac{39}{14} = \mathbf{\frac{39}{14}}$
	\end{itemize}
	
\end{document}